\documentclass{article}
\usepackage[left=2cm, right=2cm, top=2cm, bottom=2cm]{geometry}
\usepackage{graphicx}
\usepackage{amsmath}
\usepackage{amssymb}
\usepackage{amsthm}
\usepackage{fancyhdr}
\usepackage{verbatim}
\usepackage{listings}
\usepackage{xcolor}
\usepackage{pdfpages}
\usepackage{cancel}
\usepackage{physics}
\usepackage{esint}


\lstset{
    language=C++,
    basicstyle=\ttfamily\footnotesize,
    keywordstyle=\color{blue},
    commentstyle=\color{green},
    stringstyle=\color{red},
    numbers=left,
    numberstyle=\tiny\color{gray},
    frame=single,
    breaklines=true,
    captionpos=b,
}


\title{MTH 553: Lecture \# 27}
\author{Cliff Sun}

\newtheorem{theorem}{Theorem}[section]
\newtheorem{lemma}[theorem]{Lemma}
\newtheorem{definition}[theorem]{Definition}
\newtheorem{conjecture}[theorem]{Conjecture}
\newtheorem{proposition}[theorem]{Proposition}
\newtheorem{corollary}[theorem]{Corollary}
\newtheorem{one minute paper}[theorem]{One Minute Paper}

\pagestyle{fancy}
\lhead{\textbf{\thepage}\ \ \nouppercase{\rightmark}}
\chead{MTH 553: Lecture \# 27}
\rhead{Cliff Sun}

\begin{document}

\maketitle

\section*{Lecture Span}
\begin{enumerate}
    \item Mean value theorem
\end{enumerate}

\begin{proof}
    We first showed that 
    \begin{equation}
        \int_{B(x,r)}\triangle u dy = \omega_n r^{n-1}\partial_{r}\fint_{B(x,r)}
    \end{equation}
    Therefore, $u$ is harmonic if and only if the mean value is constant with respect to $r$. Thus, as $r$ goes to zero, we obtain that this converges to $u(x)$. Next, (1) implies (2) which is the
    integral over the entire ball is constant because you can just add up all $\partial B(x,r)$ from 0 to $r$. (2) implies (1) since 
    \begin{align*}
        (2) \implies u(x)\int_{B(x,r)}1 dy &= \int_{B(x,r)}u(y)dy \\
        &= u(x)\int_{\partial B(x,r)}dS = \int_{\partial B(x,r)}u dS
    \end{align*}
    This is because 
    \begin{equation}
        \int_{B} = \int_{0}^{r}\int_{\partial B}
    \end{equation}
\end{proof}

\textbf{Submean value thoerem}: Let $u \in C^{2}(\Omega)$, then if $\triangle u \geq 0$, then the mean value of $u(x)$ in a ball centered around $x$ increases as $r$ increases.  

Therefore, 

\begin{equation}
    u(x) \leq \fint_{\partial B}u dS
\end{equation}

\begin{equation}
    u(x) \leq \fint_{B}udy
\end{equation}

Remark: if $u$ is subharmonic $\triangle u \geq 0$, then $$ r^{n-1}\partial_{r}\fint_{\partial B}u dS$$ is an increasing function of $r$. This is because $$\int_B \triangle u dy$$ is increasing with respect to $r$ $\triangle u \geq 0$. 

\subsection*{Strong Maximum Principle}

Let $u$ be subharmonic in a connected open set $\Omega$. This can be possibly unbounded. Then, either $u$ is constant or $u(x)$ is strictly less then 
\begin{equation}
    u < \sup_{y \in \Omega}u(y)
\end{equation}

For all $x \in \Omega$. This means that if we find some 'max' value in the interior, then there will always be some other value of which $u(y)$ is larger. If $u \in C(\overline{\Omega})$ and $\Omega$ is bounded. Then either $u$ is constant or 

\begin{equation}
    u(x) < \max_{y \in \partial \Omega}u(y) = \max_{\overline{\Omega}}u
\end{equation}

\begin{proof}
    Here, a connected set cannot be written as a union of two disjoint sets. Assume that $A \equiv \sup_{y \in \Omega} < +\infty$. Then, we can decompose $\Omega$
    \begin{equation}
        \Omega = E \cup F
    \end{equation}  
    Where $E$ and $F$ are sets where 
    \begin{equation}
        E = \{x: \quad u(x) < A\}
    \end{equation}
    \begin{equation}
        F = \{x: \quad u(x) = A\}
    \end{equation}
    Note that $E \cap F = \empty$. $E$ is open. $F$ is also open as follows. If $x \in F$. Then $$A = u(x) \leq \frac{1}{|B(x,r)|}\int_{B}u(y)dy$$ by the sub mean value theorem. But $u \leq \sup$, thus 
    \begin{equation}
        A \leq \frac{1}{|B|}\int_{B}A dy = A
    \end{equation}
    Therefore, equality must hold all throughout our proof. Hence, 
    \begin{equation}
        u(y) = A
    \end{equation}
    For all $y \in B$. Therefore, $F$ is open because $F$ contains a neighborhood. But $\Omega$ is connected, therefore, one of the open sets $E$ or $F$ must be $\Omega$. Either $u$ is constant or $\Omega = E$. 
\end{proof}

\begin{proof}
    Either $u$ is constant or $u < A$. But $\sup_{\Omega}u = \max_{\overline{\Omega}}u = \max_{\partial \Omega}u$ since $u$ is continuous and $\overline{\Omega}$ is compact.  
\end{proof}

Consequence: continuous $L^{\infty}$ dependence on data for Dirichlet problem of Poisson's equation. If $\Omega$ is a bounded domain, then $u_1, u_2 \in C^{2} \cap C(\overline{\Omega})$. And they solve 
\begin{align*}
    -\triangle u_i &= f \\
    u_1 &= g_i
\end{align*} 

(same $f$, different $g$). With $|g_1 - g_2| < \epsilon$ on $\partial \Omega$, then $|u_1 - u_2| \leq \epsilon$ in $\Omega$.
Proof next time. 
\end{document}